% Fixed dimensions keep experiment placeholder panels visually consistent.
\newcommand{\expresultplaceholder}[1]{%
  \fbox{\begin{minipage}[c][3.35cm][c]{0.94\linewidth}%
    \centering\small\textbf{Figure placeholder}\par\vspace{3pt}%
    \textbf{#1}%
  \end{minipage}}%
}

\section{Experiments}\label{sec:experiments}

\subsection{Experimental Settings}\label{sec:exp-setup}

\noindent\textbf{Datasets.}

\textbf{T-Drive.} The dataset contains 318,774 pre-segmented Beijing taxi trajectories collected over one week. For scalability, we construct $\times1$, $\times10$, $\times50$, and $\times100$ variants by replicating trajectories with unique identifiers and shifting each replica's timestamps by one week; spatial coordinates are unchanged.

\textbf{Porto.} The dataset contains approximately 250,000 pre-segmented taxi trajectories collected over two months. Its scale variants are trajectory-level random subsamples at $\times1$, $\times0.5$, $\times0.25$, and $\times0.1$ (approximately 250,000, 125,000, 62,500, and 25,000 trajectories), generated with a fixed seed. Thus, T-Drive scaling expands the data, whereas Porto scaling measures performance on smaller samples.

\par\noindent\textbf{Query Workloads.}

\textbf{NL-to-IR.} The frozen T-Drive set contains 40 natural-language requests: 20 simple and 20 complex. It covers SRQ, ST-range, similarity search (SS), and Top-$k$, with five requests per query type and difficulty; basic and POI-oriented requests are both represented.

\textbf{Query-plan comparison.} The frozen BoundIR workloads contain 24 simple and 24 complex queries on each dataset. Each difficulty level is balanced across the same four query types (six queries per type) and is evaluated against its corresponding FullScan oracle.

\textbf{Window-distribution study.} To isolate spatial workload placement, we generate SRQ workloads under uniform, Gaussian, and skewed distributions while holding each dataset snapshot fixed. Each distribution contains ten query centers and five nested square-window widths (50, 100, 500, 1,000, and 1,500 m), with a four-hour temporal window. Uniform centers are spread over the dense spatial core; Gaussian centers cluster around its densest hotspot; skewed centers select among eight hotspots with Zipf weights (95\% of samples) and use uniform placement for the remainder. The same six query-processing methods are compared on both datasets.

\par\noindent\textbf{Baselines.}

\textbf{NL-to-IR.} DirectIR prompts a language model to emit the target IR directly. DIN-SQL (Spider/BIRD) adapts the corresponding Text-to-SQL pipelines to the query schema, while SAG-MQL uses schema-path grounding followed by MQL generation.

\textbf{Query planning and execution.} FullScan reads all raw chunks and provides the answer oracle. RBO selects access paths by a fixed priority; CBO ranks the classic access alternatives with rule-based cost estimates. Bao supplies learned plan-choice guidance, and ParamTree supplies learned cost estimates; both are adapted to the same BoundIR interface and HBase execution path.

\par\noindent\textbf{Metrics and Parameter Settings.}
We evaluate translation action accuracy and query-processing latency, with correctness checked against FullScan. Table~\ref{tab:exp-settings} summarizes the beam and MC settings and the supported query types.

\FloatBarrier
\begin{table}[tbp]
\caption{Query and search parameter settings.}\label{tab:exp-settings}
\centering\small
\renewcommand{\arraystretch}{1.12}
\begin{tabularx}{0.62\textwidth}{@{}P{0.27\textwidth}Y@{}}
\toprule
Parameter & Setting\\
\midrule
Beam & $\{1,3,5,7\}$ (default: 3)\\
MC & $\{4,6,8,10\}$ (default: 4)\\
Query type & SRQ, ST-range, SS, Top-$k$\\
\bottomrule
\end{tabularx}
\end{table}

\par\noindent\textbf{Implementation Details.}
All plan and execution comparisons use frozen BoundIR inputs rather than generated NL-to-IR outputs, so translation errors do not affect optimizer measurements. For each query and dataset, every method receives the same snapshot and query; its result is compared with the FullScan oracle (set equality for identifier queries and ordered-ID equality for Top-$k$). The E1 translation suite is run separately against its expected action and gold IR. Runs use HBase 2.2.3, Hadoop 2.8.5, ZooKeeper 3.4.10, OpenJDK 8, and four RegionServers. The advanced indexes are enabled for workloads that use them; caches are warmed before timed repetitions.



\subsection{Overall Performance}\label{sec:exp-overall}

We evaluate translation quality, workload sensitivity, and query planning and execution. The NL-to-IR results summarize E1, the distribution table characterizes SRQ workloads, and the E2/E3 plots compare planning and end-to-end performance across query types, difficulty levels, and datasets.

\noindent\textbf{Workload results.}
Table~\ref{tab:exp-workload-results} reports end-to-end latency for the uniform, Gaussian, and skewed SRQ workloads. Each cell is the P50 latency in milliseconds; entries without recorded measurements are marked \TBD.

\begin{table}[tbp]
\caption{End-to-end latency (P50, ms) by spatial workload distribution.}\label{tab:exp-workload-results}
\centering\small
\renewcommand{\arraystretch}{1.08}
\begin{tabularx}{\textwidth}{@{}P{0.12\textwidth}P{0.15\textwidth}*{6}{>{\centering\arraybackslash}X}@{}}
\toprule
Dataset & Distribution & \KART{} & FullScan & RBO & CBO & Bao & ParamTree\\
\midrule
\multirow{3}{*}{\parbox{0.12\textwidth}{\centering T-Drive}} & Uniform & \TBD & \TBD & \TBD & \TBD & \TBD & \TBD\\
 & Gaussian & \TBD & \TBD & \TBD & \TBD & \TBD & \TBD\\
 & Skewed & \TBD & \TBD & \TBD & \TBD & \TBD & \TBD\\
\addlinespace[2pt]
\multirow{3}{*}{\parbox{0.12\textwidth}{\centering Porto}} & Uniform & \TBD & \TBD & \TBD & \TBD & \TBD & \TBD\\
 & Gaussian & \TBD & \TBD & \TBD & \TBD & \TBD & \TBD\\
 & Skewed & \TBD & \TBD & \TBD & \TBD & \TBD & \TBD\\
\bottomrule
\end{tabularx}
\end{table}

\par\noindent\textbf{NL-to-IR (E1).}
Table~\ref{tab:exp-nl-results} presents the expected-action accuracy and P50 parse latency for simple and complex requests. Action accuracy is shown separately for Pass, Reject, and Clarify; Pass additionally requires semantic equivalence to the gold IR.

\begin{table}[tbp]
\caption{E1 NL-to-IR results on T-Drive. Action columns report Pass/Reject/Clarify accuracy; latency is P50 in milliseconds.}\label{tab:exp-nl-results}
\centering\small
\renewcommand{\arraystretch}{1.15}
\begin{tabularx}{\textwidth}{@{}P{0.20\textwidth}*{4}{>{\centering\arraybackslash}X}@{}}
\toprule
 & \multicolumn{2}{c}{Simple queries} & \multicolumn{2}{c}{Complex queries}\\
\cmidrule(lr){2-3}\cmidrule(lr){4-5}
Method & \shortstack{Action\\accuracy} & P50 & \shortstack{Action\\accuracy} & P50\\
\midrule
\KART{} & 1.00/1.00/1.00 & 1,562 & 1.00/1.00/1.00 & 1,596\\
DirectIR & 0.7778/0.3333/0 & 1,586 & 0.7778/0.3333/0 & 2,033\\
DIN-SQL (Spider) & 0.3333/0.3333/0 & 14,151 & 0.6667/0.5000/0 & 14,097\\
DIN-SQL (BIRD) & 0.6667/0.3333/0 & 18,292 & 0.5556/0.5000/0 & 8,272\\
SAG-MQL & 0.6667/0.5000/0 & 1,329 & 0.7778/0.3333/0 & 1,789\\
\bottomrule
\end{tabularx}
\end{table}

\FloatBarrier
\par\noindent\textbf{Query planning and execution (E2/E3).}
E2 measures planning latency from frozen BoundIR inputs; E3 measures planning, validation, and physical execution end to end. The six panels in Fig.~\ref{fig:exp-overall-e2e} compare the methods by dataset and report planning and execution behavior without aggregating the simple and complex workloads together. All query outputs are checked against FullScan.

\begin{figure}[!htbp]
\centering
\begin{minipage}[t]{0.31\textwidth}\centering
\expresultplaceholder{(a) E2 planning: T-Drive}
\end{minipage}\hfill
\begin{minipage}[t]{0.31\textwidth}\centering
\expresultplaceholder{(b) E2 planning: Porto}
\end{minipage}\hfill
\begin{minipage}[t]{0.31\textwidth}\centering
\expresultplaceholder{(c) E3 end-to-end: T-Drive}
\end{minipage}\\[4pt]
\begin{minipage}[t]{0.31\textwidth}\centering
\expresultplaceholder{(d) E3 end-to-end: Porto}
\end{minipage}\hfill
\begin{minipage}[t]{0.31\textwidth}\centering
\expresultplaceholder{(e) E3 latency breakdown: T-Drive}
\end{minipage}\hfill
\begin{minipage}[t]{0.31\textwidth}\centering
\expresultplaceholder{(f) E3 latency breakdown: Porto}
\end{minipage}
\caption{Planning and end-to-end performance across datasets and workload difficulty.}\label{fig:exp-overall-e2e}
\end{figure}
\FloatBarrier

\subsection{Ablation Study}\label{sec:exp-ablation}

The ablation compares full \KART{} with variants that remove the rewrite-rule LLM, the cost-model LLM, the cost model, or plan verification.
  The results show exact-match agreement of 1.00 for all tested variants. Removing the cost model causes the largest E2E increase, particularly on Porto complex queries; the other removals have smaller, workload-dependent effects.

\begin{figure}[!htbp]
\centering
\begin{minipage}[t]{0.48\textwidth}\centering
\expresultplaceholder{(a) Ablation E2E latency}
\end{minipage}\hfill
\begin{minipage}[t]{0.48\textwidth}\centering
\expresultplaceholder{(b) E2E increase over Full \KART{}}
\end{minipage}
\caption{Ablation results across datasets and workload difficulty.}\label{fig:exp-ablation}
\end{figure}
\FloatBarrier

\noindent\textbf{Rewrite-rule LLM.} Removing conditional rewrite rules increases E2E latency most clearly on T-Drive, where cost-guided candidates lose useful access-path refinements. The smaller Porto change suggests that the benefit depends on how well the recorded rules match a dataset's query and layout patterns.

\noindent\textbf{Cost-model LLM.} Without offline cost calibration, plan ranking becomes less aligned with observed execution cost. The increase is modest on T-Drive but pronounced for complex Porto queries, where selecting an expensive plan can dominate E2E latency.

\noindent\textbf{Cost model.} Removing cost estimates has the largest impact. Fixed or weakly differentiated choices cannot account for scan, candidate-fetch, and refinement costs, so the selected access path can perform substantially more execution work.

\noindent\textbf{Plan verification.} Disabling verification does not consistently lower or raise E2E latency, and the tested outputs remain exact. Verification instead enforces semantic, coverage, and physical-safety conditions that timing results alone do not establish.

% Parameter, scalability, and case-study experiments included in the main paper.

\subsection{Parameter Study}\label{sec:exp-parameters}

We study two search-budget parameters. Beam width (beam) limits the number of partial plans retained at each search step, while the candidate limit (MC) bounds the number of completed plans considered. We report E2E latency only. In the beam sweep, MC is fixed at 8; in the MC sweep, beam is fixed at 3. The configured defaults are beam 3 and MC 4.

\begin{figure}[!htbp]
\centering
\begin{minipage}[t]{0.31\textwidth}\centering
\expresultplaceholder{(a) Beam width}
\end{minipage}\hfill
\begin{minipage}[t]{0.31\textwidth}\centering
\expresultplaceholder{(b) Candidate limit (MC)}
\end{minipage}\hfill
\begin{minipage}[t]{0.31\textwidth}\centering
\expresultplaceholder{(c) Cross-dataset summary}
\end{minipage}
\caption{E2E latency under different bounded-search parameter settings.}\label{fig:exp-parameters}
\end{figure}
\FloatBarrier

For beam width, T-Drive reaches its lowest recorded E2E latency at beam 3 (195/220 ms for simple/complex queries); wider beams do not yield consistent gains (199--205/232--237 ms). Porto is less sensitive, with recorded values between 1,043 and 1,060 ms for simple queries and 1,043 and 1,057 ms for complex queries. This suggests that retaining a wider frontier adds search effort without reliably improving the selected plan.

For MC, T-Drive improves as the limit grows from 4 to 8 (207/239 to 195/220 ms), then regresses at 10 (206/244 ms), indicating a useful balance between considering more alternatives and limiting planning overhead. Porto varies less monotonically: the simple workload is lowest at MC 10 (1,028 ms), while complex latency is similar at MC 4 and 10 (1,043 ms) and slightly higher at intermediate settings. Thus, the best MC value depends on dataset and workload difficulty; the recorded changes are in E2E latency only.

\subsection{Scalability}\label{sec:exp-scalability}

We measure E2E latency across T-Drive expansions ($\times1$, $\times10$, $\times50$, $\times100$) and Porto trajectory subsamples ($\times1$, $\times0.5$, $\times0.25$, $\times0.1$). Each scale uses the same query semantics and compares \KART{} with the five IR-level baselines. The measured latency and index-footprint results are \TBD.

\begin{figure}[!htbp]
\centering
\begin{minipage}[t]{0.31\textwidth}\centering
\expresultplaceholder{(a) T-Drive expansion}
\end{minipage}\hfill
\begin{minipage}[t]{0.31\textwidth}\centering
\expresultplaceholder{(b) Porto subsampling}
\end{minipage}\hfill
\begin{minipage}[t]{0.31\textwidth}\centering
\expresultplaceholder{(c) Data and index scale}
\end{minipage}
\caption{Query latency and data/index size across dataset scales.}\label{fig:exp-scalability}
\end{figure}
\FloatBarrier

The expected T-Drive trend is increasing E2E latency with larger replicas, since more rows and postings increase index scans and chunk retrieval. With fixed query windows and spatial coordinates, however, the candidate set need not grow in proportion to the replicated dataset; indexed plans should therefore scale more gently than FullScan. For Porto, smaller samples reduce the indexed population, candidate chunks, and exact-refinement work, so latency is expected to fall as the sample fraction decreases. At the smallest scales, fixed planning and scan startup costs may limit further gains. The measured trends are \TBD; these effects explain the expected dataset-specific behavior.

\subsection{Representative Query Case Study}\label{sec:exp-case-study}

\TBD

\begin{figure}[!htbp]
\centering
\expresultplaceholder{Representative query case study}
\caption{Case-study query, selected access plan, and result.}\label{fig:exp-case-study}
\end{figure}
\FloatBarrier


